% Theorem statements for "The Price of Self-Knowledge" (ML x OR @ NeurIPS 2026).
% Compiles standalone (pdflatex theorems.tex); for the NeurIPS build, \input
% the body between the BEGIN/END BODY markers and keep the macros.
% Numbering here matches THEOREMS.md; renumber only via that register.
\documentclass[11pt]{article}
\usepackage[margin=1.1in]{geometry}
\usepackage{macros}  % shared with proofs.tex and the NeurIPS main.tex

\begin{document}
\section*{Theorem statements (register copy)}

% ================= BEGIN BODY =================

\begin{lemma}[Cost equivalence; L1]
Under A1--A3 with conditionally symmetric exploration, the expected
incremental performative cost satisfies
$C_T = \tfrac{1}{2}\gpo\,\E\!\big[\sum_{t\le T}\dev_t^2\big] + R_T$ with
$|R_T| \le \tfrac{L_3}{6}\,\E\sum_t |\dev_t|^3$, and the realized cost's
variance decomposes exactly as
$\Var(C_T^{\mathrm{real}}) = \delta_0^2\,\Var\!\big(\sum_t \dev_t\big)
+ \tfrac{\gpo^2}{4}\,\Var\!\big(\sum_t \dev_t^2\big)$
with vanishing cross term.
\end{lemma}

\begin{theorem}[Information saturation; T1]
With no exploration and $\rho(M)<1$, the total design energy of the
retraining trajectory equals $\dev_0^\top P\,\dev_0$, where $P$ is the
unique solution of the discrete Lyapunov equation $P = I + M^\top P M$.
Consequently the Fisher information for the response Jacobian is bounded,
in every direction, uniformly in the horizon:
$\Var(\hat\eps) \ge \sigma^2\big(\dev_0^\top P_u\,\dev_0\big)^{-1}$ for the
component along any direction $u$.
\end{theorem}

\begin{corollary}[Safety implies blindness; C1.1]
In the scalar case the information cap $(h_0-\hstar)^2/(1-m^2)$ is strictly
increasing in the modulus $m$: a fast-contracting loop reveals least about
its own performativity, and precise self-knowledge is available only near
the stability boundary.
\end{corollary}

\begin{corollary}[Trajectory design degeneracy; C1.2]
The noiseless trajectory's empirical design has at most two effective
support points (geometric collapse for $m<1$; a period-two orbit at
$m=1$).
\end{corollary}

\begin{remark}[Non-normal transient information; R1]
The energy identity $\dev_0^\top P\,\dev_0$ of T1 holds for non-normal $M$
as well, where the energy can exceed the spectral cap
$\|\dev_0\|^2/(1-\rho(M)^2)$ by up to the factor $\kappa(V)^2$.
\end{remark}

\begin{theorem}[The exchange rate, pathwise; T2]
In the stationary regime, on every realization,
\[
\Var(\hat\eps \mid \mathrm{design}) \times \tfrac{1}{2}\gpo \sum_{t\le T}
\dev_t^2 \;=\; \tfrac{1}{2}\,\gpo\,\sigma^2\Big(1 +
\tfrac{T\,\bar{\dev}^2}{\Sxx}\Big)
\;=\; \tfrac{1}{2}\,\gpo\,\sigma^2\,\big(1 + O_P(1/T)\big),
\]
independent of the exploration intensity, the horizon, and the modulus.
\end{theorem}

\begin{proposition}[Concentration; P2.1]
In the stationary regime the measured product concentrates around the
invariant at rate $T^{-1/2}$, with long-run
$\Var\big(\sum_{t\le T}\dev_t\big) = T\,v\,\tfrac{1-m}{1+m}\,(1+o(1))$.
\end{proposition}

\begin{proposition}[Feedback bias; P2.2]
Under the feedback relaxation A3$'$ with gain $\phi$, the OLS response
slope carries the bias
\[
\E[\hat\eps - \eps_{\mathrm{true}}] \;=\;
\frac{\phi\,\sigma^2\,(3m-1)}{(1-m^2)\,T\,v_\phi}\,(1+o(1)),
\qquad v_\phi = \frac{\sigma_e^2 + \phi^2\sigma^2}{1-m^2},
\]
which is $O(1/T)$ and changes sign at $m = 1/3$.
\end{proposition}

\begin{theorem}[Minimax lower bound, deviation budget; T3]
Over all non-anticipating exploration policies with
$\sum_{t\le T} \dev_t^2 \le D$ almost surely and all estimators (biased
included), the minimax risk over an interval of width $w$ satisfies
\[
\inf_{\text{policy, estimator}}\;\sup_{\eps}\;
\E\big[(\hat\eps - \eps)^2\big] \;\ge\;
\frac{\sigma^2}{D + \sigma^2 (\pi/w)^2},
\]
so the exchange-rate product is a floor, attained with equality by
symmetric probing with least squares.
\end{theorem}

\begin{lemma}[Exploitation--information; L2]
For the symmetric two-point response family with half-width $\delta$ and
any non-anticipating policy under the certainty-equivalent anchor,
\[
\E[\mathrm{gain}_T] \;\le\;
\frac{g_1\,\delta^{3/2}\,T^{1/2}\,\big(\E S_T\big)^{3/4}}{\sigma^{1/2}},
\qquad g_1 = \beta\,|\hstar - \psi|,
\]
via the exact identity $\E\,\big|\,\E[s\mid\mathcal{F}_t]\,\big|
= \TV(P^+_t, P^-_t)$ and Pinsker's inequality
$\TV_t \le \delta\sqrt{\E S_t}/\sigma$: exploitation is bounded by the
information already purchased. At the minimax scale
$\delta = c\,\sigma/\sqrt{\E S_T}$ this is $O(\sigma\sqrt{T})$, an
$O(T^{-1/2})$ fraction of the $\Theta(T)$ exploration cost.
\end{lemma}

\begin{theorem}[Minimax lower bound, value budget; T4]
Under stationary-scale exploration and the certainty-equivalent anchor,
for every non-anticipating policy and every estimator,
\[
\sup\;\;\Var(\hat\eps)\times C_T \;\ge\; \frac{\gpo\,\sigma^2}{27}\,
\big(1 - O(T^{-1/2})\big),
\]
the supremum over the two-point response family at the minimax scale and
the constant arising from the Le Cam reduction optimized at $c = 2/3$.
The exploited rebate vanishes at rate $T^{-1/2}$ (Lemma L2). The sharp
constant $\tfrac{1}{2}$, against which no simulated policy fell by more
than Monte-Carlo error (\texttt{run\_open1.py} \S B), is conjectured and
is part of OPEN-1.
\end{theorem}

\begin{theorem}[Optimal exploration shapes; T5]
Under the value budget $\tr(\Gpo M) \le B$ on the exploration second
moment $M$:
(a) the A-optimal design is $M^{*} = \big(B/\tr\,\Gpo^{1/2}\big)\,
\Gpo^{-1/2}$ with value $(\tr\,\Gpo^{1/2})^2/B$;
(b) the D-optimal design is $M^{*} = (B/d)\,\Gpo^{-1}$;
(c) the c-optimal design for the correction functional $c$ is
$M^{*} = B\,cc^\top/(c^\top \Gpo\, c)$ (probing along the correction
direction itself), with value $c^\top \Gpo\, c / B$.
\end{theorem}

\begin{corollary}[Price of isotropic jitter; C5.1]
Isotropic exploration at matched budget overpays the A-optimal design by
exactly the curvature-dispersion factor
$F = d\,\tr(\Gpo)/(\tr\,\Gpo^{1/2})^2 \ge 1$, with equality iff all
curvature directions are equal.
\end{corollary}

\begin{lemma}[Temporal shaping is irrelevant; L3]
With serially independent response noise, the information about the
response parameters depends on the exploration process only through its
empirical second moment; temporally correlated exploration buys nothing
beyond its stationary design measure.
\end{lemma}

\begin{theorem}[Chebyshev unidentifiability; T6]
The sensitivity system of the structural family
$\tau(h) = C_0 + C_1 e^{-ch}$ is an extended Chebyshev system; any design
with fewer than three support points has singular Fisher information, and
by C1.2 no retraining trajectory, at any modulus, identifies the response
curvature $c$.
\end{theorem}

\begin{lemma}[Perturbed modulus; L4]
If the corrected update uses an estimated response with $C^1$ error $e$,
the closed-loop modulus satisfies
$|\hat\rho - \rho| \le \eta\,\beta\,\big(\|e\|_\infty +
|\hstar-\psi|\,\|e'\|_\infty\big)$.
\end{lemma}

\begin{remark}[Open-loop neutrality; R2]
History-independent exploration schedules cannot change the linearized
contraction rate of the retraining loop.
\end{remark}

\begin{proposition}[Safety under pessimism; P7.1]
Given an anytime-valid confidence sequence for the response parameters at
level $1-\delta$, the pessimism rule (worst-case modulus over the
confidence set kept below $1 - c_{\mathrm{margin}}$, correction frozen
otherwise) keeps the realized modulus within the margin for all $t$
simultaneously with probability at least $1-\delta$.
\end{proposition}

\begin{theorem}[Anchoring crossover, horizon-dependent; T7]
Under the value budget $B$ and horizon $T$, the budget-and-horizon-optimal
nonparametric (secant) estimator has
\[
\mathrm{MSE}_{\mathrm{np}} = \frac{\gpo\,\sigma^2}{2B}
+ \Big(\frac{\tau'''(\hstar)}{6}\Big)^{\!2}
\Big(\frac{2B}{\gpo\,T}\Big)^{\!2},
\]
so anchoring to a structural family with misspecification
$\delta_{\mathrm{mis}}$ wins iff (to leading order)
$\delta_{\mathrm{mis}} < |\tau'''(\hstar)|\,B/(3\,\gpo\,T)$: anchoring wins
at short horizons and tight budgets, the free-form route asymptotically in
$T$.
\end{theorem}

\begin{theorem}[The ROI of self-knowledge; T8]
Exploring once to accuracy $v$ and correcting forever, discounted at
$\rho_{\mathrm{disc}}$, is optimal at
$v^{*} = (\sigma/\kappa)\sqrt{\rho_{\mathrm{disc}}}$, and exploring at all
is worthwhile iff
\[
\rho_{\mathrm{disc}} \;<\; \rho^{*} \;=\;
\frac{(\hsp - \hpo)^4}{4\,\kappa^2\,\sigma^2},
\]
in which the curvature $\gpo$ cancels completely.
\end{theorem}

\begin{proposition}[Frontier consistency; P9.1]
Every deterministic design sits exactly on the exchange-rate frontier; in
particular the $O(\log T)$-cost adaptive-design schedules of Lai--Robbins.
\end{proposition}

\begin{proposition}[Isotropy contrast; P9.2]
Naive isotropic exploration, optimal in isotropic LQR-like settings,
overpays by exactly the dispersion factor $F$ in the performative
setting, with $F = 1$ iff the curvature spectrum is flat.
\end{proposition}

\begin{proposition}[Separable multi-bond curvature; P9.3]
In the separable multi-bond market,
$(\Gpo)_{aa} = \gamma_a + \beta\,\eps_a\,(2 + c_t\psi_a - c_t h_a)$
exactly, with the factor-coupled case holding to first order in the
off-diagonal response.
\end{proposition}

% ================= END BODY =================
\end{document}
